\documentclass[11pt,letterpaper]{article}
\usepackage[T1]{fontenc}
\usepackage[utf8]{inputenc}
\usepackage{lmodern,microtype}
\usepackage[margin=1in]{geometry}
\usepackage{amsmath,amssymb,amsthm,mathtools}
\usepackage{booktabs,array,longtable}
\usepackage{xcolor}
\usepackage{enumitem}
\usepackage{listings}
\usepackage{hyperref}
\hypersetup{colorlinks=true,linkcolor=blue!45!black,citecolor=blue!45!black,urlcolor=blue!45!black,
 pdftitle={Component Compression and Restart Entropy in Exact Unknot Recognition},
 pdfauthor={Research contribution prepared for ProveIt}}
\setlength{\parskip}{0.35em}
\setlength{\parindent}{1.2em}
\setlength{\emergencystretch}{2em}
\allowdisplaybreaks
\newtheorem{theorem}{Theorem}[section]
\newtheorem{proposition}[theorem]{Proposition}
\newtheorem{lemma}[theorem]{Lemma}
\newtheorem{corollary}[theorem]{Corollary}
\theoremstyle{definition}
\newtheorem{definition}[theorem]{Definition}
\newtheorem{example}[theorem]{Example}
\newtheorem{question}[theorem]{Research question}
\theoremstyle{remark}
\newtheorem{remark}[theorem]{Remark}
\newcommand{\F}{\mathbb F_2}
\newcommand{\R}{\mathcal R}
\newcommand{\C}{\mathcal C}
\newcommand{\supp}{\operatorname{supp}}
\newcommand{\rank}{\operatorname{rank}}
\newcommand{\Kh}{\operatorname{Kh}}
\newcommand{\poly}{\operatorname{poly}}
\newcommand{\Comp}{\operatorname{Comp}}
\newcommand{\eps}{\varepsilon}
\newcommand{\bits}{\operatorname{bits}}
\newcommand{\ceil}[1]{\left\lceil #1\right\rceil}
\newcommand{\floor}[1]{\left\lfloor #1\right\rfloor}
\lstset{basicstyle=\ttfamily\small,columns=fullflexible,keepspaces=true,
 breaklines=true,frame=single,framerule=0.3pt,showstringspaces=false,
 xleftmargin=0.4em,xrightmargin=0.4em}
\title{\textbf{Component Compression and Restart Entropy\\in Exact Unknot Recognition}\\[0.6em]
\large Arithmetic improvements and a conditional route to quasi-polynomial complexity}
\author{Research contribution prepared for the ProveIt project}
\date{7 October 2026}
\begin{document}
\maketitle
\begin{abstract}
We develop two complementary improvements to ProveIt's unknot-recognition
program. First, a compiled genus-zero dotted-cobordism composition factors through
a square-zero algebra with one variable per connected gluing component. Ranked
subset convolution then replaces dense pairs of input monomials by a
$\poly(w)2^{w/2}$-bit-time composition algorithm at boundary size $w$.
The component quotient is the exact observational quotient for all opposite
operands; in particular, it is dimension-optimal among universal linear
summaries. An explicit family of planar matching triples realizes arbitrarily
large input circle spaces with a fixed number of gluing components.
A standard-library Python implementation is cross-checked against an excerpt
of the repository's numerical routines, with 460,564 matching-composition
comparisons and 1,000 associativity checks. Paired dense numerical-kernel
benchmarks show substantial improvements, while a deliberately sparse case
shows why unconditional replacement would be inappropriate.

Second, we refine the combinatorics of lexicographic hierarchy restarts.
If length-$L$ complexity profiles have coordinates at most $g$ and at most
$s$ positive coordinates, the exact state bound is
$\sum_{j\leq s}\binom Lj g^j$, rather than $(g+1)^L$.
This yields quasi-polynomial restart counts when $L,g$ are polynomial and
$s=O(\log n)$, even without logarithmic total depth. We give an exact rank
potential, a sharpness construction, and executable finite-trace checks.
The topology needed to establish this sparsity uniformly, and the complete
bit-cost bounds for hierarchy construction, remain explicit research
obligations. This work does not establish an unrestricted quasi-polynomial
unknot recognizer.
\end{abstract}
\vspace{0.3em}
\noindent\textbf{Artifact status.} Complete proofs of the algebraic and counting
statements; tested numerical kernel and trace tools; opt-in integration harness.
No full-recognizer speedup, full-repository test run, or first-in-literature
priority is asserted.
\newpage
\begingroup
\setlength{\parskip}{0pt}
\small
\tableofcontents
\endgroup
\newpage

\section{Scope, baseline, and the two complexity problems}
\label{sec:scope}
The practical objective is a reliable decision procedure for a knot diagram,
with an eventual target of $n^{O(\log n)}$ running time. The relevant ProveIt
subtree contains an exact implementation and a separate, unfinished
hierarchy-based complexity program. These should be improved together, but
not conflated.

The audited Python pipeline first applies diagram simplification and
one-sided invariant filters, then uses a scanning Khovanov computation when
necessary. Its numerical backend already interns matchings, compiles gluing
plans, caches shape-dependent results, and cancels units with a sparse
elimination strategy. Visible connected-sum factorization and multiple scan
orders are also existing features. The repository describes the complete
fallback as exponential, not quasi-polynomial \cite{repo}.

This article targets a remaining dense-arithmetic cost in
\texttt{fastunknot/planar.py}: on a numerical composition-cache miss, the
implementation iterates over the support of each operand and evaluates every
pair. The exact audited blob is
\texttt{1078526e7e7dbaf0b7267728105d870d94136769}.
The scanner interface was checked against blob
\texttt{2c1ad52d14296b109376af19668ae0eebb4c6fdd} of
\texttt{scan\_fast.py}. These are content identifiers, not a claim to have
executed a complete checkout.

There are two different sources of large computations. One is the cost of
composing two already represented morphisms. The other is the number of
objects and morphisms that the computation must represent. Improving the
first does not automatically control the second. Our component-compression
theorem addresses the first; the multiplicity example in
Section~\ref{sec:multiplicity} identifies a genuine obstruction to using it
as a universal width-only theorem.

The hierarchy direction has a different distinction: a bound on the number
of restarts is not yet a bound on the cost of finding surfaces, cutting along
them, and maintaining their descriptions. We formulate the sparse-profile
criterion so that these additional obligations cannot disappear inside an
unspecified primitive.

\subsection{Contributions and their status}
The main unconditional contribution is an exact factorization of the
repository's composition law through its connected gluing components,
together with a dense-bit implementation and its correctness proof.
We prove that this quotient retains precisely the information that any
universal opposite-operand composition query can observe. We also give
explicit planar ``zipper'' triples, rather than relying only on artificial
component arrays, to exhibit large compression ratios.

The second unconditional contribution is a sparse lexicographic restart
potential. Its count is elementary, but its significance for the research
program is concrete: one may seek logarithmically many positive-complexity
layers instead of logarithmic total hierarchy length. The conditional
algorithmic theorem spells out the constructive topological assumptions
needed to exploit that change.

Fast subset convolution itself is not new; our implementation uses the
ranked zeta-transform method of Bj\"orklund, Husfeldt, Kaski, and Koivisto
\cite{bhkk}. Dotted-cobordism and local tangle simplification are likewise
standard foundations \cite{bn-tangles,bn-fast}. The contribution claimed
here is the proved application, the exact quotient analysis, the concrete
implementation, and the refined research target. A literature-priority
claim would require a broader comparison than the targeted source review
performed for this article.

\section{Dotted-cobordism conventions}
\label{sec:conventions}
All coefficients are in $\F$, and all vector spaces in the numerical
analysis forget quantum gradings. Homological degrees remain part of the
scanner's objects. A universal statement about arbitrary coefficient vectors
below should not be silently restricted to one graded subspace.

For a finite set $I$, put
\[
  \R(I)=\F[x_i:i\in I]/(x_i^2:i\in I),
  \qquad x^S=\prod_{i\in S}x_i.
\]
Write $\R_r=\R([r])$, using either zero-based or one-based labels as long as
the choice is consistent. Its vector-space dimension is $2^r$.
A morphism on $r$ circles is represented by an integer
\[
  \operatorname{enc}(f)=\sum_{S\subseteq[r]} f_S\,2^{\operatorname{mask}(S)}.
\]
Thus a subset mask and the position of its coefficient bit are distinct
levels of encoding. Addition of morphisms is bitwise XOR. An integer may
contain $2^r$ coefficient bits; operations on it are not unit-cost operations.

The Frobenius algebra on one circle is $A=\F[x]/(x^2)$, with
\[
 \eps(1)=0,\quad \eps(x)=1,\qquad
 \Delta(1)=1\otimes x+x\otimes1,\quad
 \Delta(x)=x\otimes x.
\]
These are the usual characteristic-two dotted-cobordism rules
\cite{bn-tangles}. The handle operator $m\Delta$ vanishes: it sends $1$ to
$2x=0$ and $x$ to $x^2=0$. Consequently a connected component of positive
genus contributes zero.

\begin{definition}[Compiled component plan]
A genus-zero plan is a list
\[
 (L_j,R_j,B_j,e_j),\qquad 1\leq j\leq k,
\]
where $L_j$, $R_j$, and $B_j$ partition the left input circles, right input
circles, and output circles respectively. Empty parts are allowed in the
general algebraic statement. The nonnegative integer $e_j$ is the number of
extra dots. The three circle universes have sizes $r,s,t$.
\end{definition}

For one pair of input monomials $x^U,z^V$, component $j$ carries
\[
 d_j=|U\cap L_j|+|V\cap R_j|+e_j
\]
dots. If $d_j\geq2$, the answer is zero. For $d_j=1$, it contributes the
monomial with all output circles in $B_j$ dotted. For $d_j=0$ and
$B_j\ne\varnothing$, it contributes
\(\sum_{b\in B_j}x^{B_j\setminus\{b\}}\).
For $B_j=\varnothing$ the contribution is the scalar $\eps(x^{d_j})$.
The full value is the product of component contributions, extended
bilinearly in the inputs. This is precisely the numerical rule used by the
source-excerpt reference harness.

A component in \texttt{planar.py} has a fifth, redundant field containing
the possible undotted-output choices. The new implementation accepts either
four or five fields and checks the fifth field for consistency. It also
checks disjointness of masks and contiguous circle indices. Geometry and
circle-numbering conventions are not recomputed by the replacement kernel.

\section{Factorization through connected components}
\label{sec:factor}
Let $y_1,\ldots,y_k$ be variables indexed by the gluing components. Define
algebra homomorphisms
\[
 P_L:\R_r\longrightarrow\R_k,
 \quad x_i\longmapsto y_j\quad(i\in L_j),
 \qquad
 P_R:\R_s\longrightarrow\R_k,
 \quad z_i\longmapsto y_j\quad(i\in R_j).
\]
These maps are well defined because the image of each square is zero.
Different input variables on one component acquire the same image, so a
monomial containing two such variables vanishes. Surviving monomials with
the same component subset are added modulo two.

For a finite output block $B$ define a linear map
$D_B:A\to\R(B)$ by
\[
\begin{array}{lll}
 B\ne\varnothing:&D_B(1)=\displaystyle\sum_{b\in B}x^{B\setminus\{b\}},
                  &D_B(x)=x^B,\\[0.5em]
 B=\varnothing:&D_B(1)=0,&D_B(x)=1.
\end{array}
\]
Set $\Omega=\bigotimes_{j=1}^k D_{B_j}$, with tensor factors ordered by the
compiled component list. This is a linear reconstruction map, not generally
an algebra homomorphism.

\begin{theorem}[Component factorization]
\label{thm:factor}
If every $e_j$ is zero or one and
$E=\{j:e_j=1\}$, then the compiled composition satisfies
\begin{equation}
 \Comp(f,g)=\Omega\!\left(y^E P_L(f)P_R(g)\right).
 \label{eq:factor}
\end{equation}
If some $e_j\geq2$, the composition is identically zero. A positive-genus
component also makes the composition identically zero over $\F$.
\end{theorem}
\begin{proof}
It suffices to take monomial inputs. If two left dots lie on one component,
$P_L$ sends the monomial to zero; the same holds on the right. If one left
and one right dot lie on the same component, their product contains $y_j^2$
and is zero. Multiplication by $y^E$ imposes exactly the analogous condition
for extra dots. Thus the product survives precisely when every component
has at most one dot, and its $j$th exponent is then the total number of dots
on that component. Applying $D_{B_j}$ gives its prescribed component value,
including the counit when the output block is empty. Multiplying these
values gives the reference rule. Bilinearity proves the identity for all
inputs. The two vanishing statements follow from $x^2=0$ and $m\Delta=0$.
\end{proof}

The quotient calculation is stronger than merely rejecting incompatible
monomial pairs early. Terms can cancel \emph{before} the two operands are
multiplied. For example, two single dots on different circles of the same
component have equal images and their sum disappears.

\subsection{Reconstruction without branching over outputs}
\begin{lemma}[Unique component preimage of an output coefficient]
\label{lem:output}
Fix an output subset $X\subseteq[t]$. For each component with a nonempty
output block, count the missing dots $|B_j\setminus X|$. If this count is
at least two for any $j$, the coefficient of $x^X$ in $\Omega(h)$ is zero.
Otherwise define $Z(X)\subseteq[k]$ by including $j$ exactly when
$B_j=\varnothing$ or $B_j\subseteq X$. Then
\[
 [x^X]\Omega(h)=[y^{Z(X)}]h.
\]
In particular, in \eqref{eq:factor} the output coefficient is zero unless
$E\subseteq Z(X)$; when this condition holds it is
$[y^{Z(X)\setminus E}]P_L(f)P_R(g)$.
\end{lemma}
\begin{proof}
For a nonempty block, $D_B(x)$ has every output circle dotted, whereas each
monomial of $D_B(1)$ has exactly one missing dot. These supports are disjoint.
For an empty block, only $D_B(x)$ survives. The component input exponent is
therefore uniquely determined by $X$. Tensoring the blocks gives the first
formula. Multiplication by $y^E$ kills monomials already meeting $E$ and
adjoins $E$ to every surviving monomial, giving the second formula.
\end{proof}

This lemma permits a dense pass over $2^t$ output coefficients. It avoids
materializing a separate expansion of $\Omega(y^S)$ for every component
subset $S$.

\begin{example}[Cancellation in a two-component endomorphism algebra]
Let every component have one left, one right, and one output circle. Then
$P_L,P_R,\Omega$ are the identity identifications. For
\[
 f=1+x_1+x_1x_2,\qquad g=1+x_2,
\]
one obtains $fg=1+x_1+x_2$, since the two $x_1x_2$ terms cancel. With
zero-based subset masks, the corresponding integer computation is
$\Comp(11,5)=7$. This example is included in
\texttt{results/examples.json}.
\end{example}

\section{Fast exact arithmetic and bit complexity}
\label{sec:arithmetic}
\subsection{Disjoint subset convolution}
Write $a=P_L(f)$ and $b=P_R(g)$. Multiplication in $\R_k$ is the disjoint
subset convolution
\begin{equation}
 h(S)=\sum_{T\subseteq S} a(T)b(S\setminus T),\qquad S\subseteq[k].
 \label{eq:conv}
\end{equation}
The direct implementation uses $3^k$ scalar products. More importantly for
the repository, multiplying before projection can use as many as $2^{r+s}$
input-monomial pairs. Projection must therefore precede convolution.

For completeness, we give the ranked-transform argument underlying the
classical $O(k^2 2^k)$ ring-operation bound \cite{bhkk}. Define
\[
 a_i(S)=\sum_{\substack{T\subseteq S\\|T|=i}}a(T),\qquad
 b_i(S)=\sum_{\substack{T\subseteq S\\|T|=i}}b(T),
 \qquad c_\ell(S)=\sum_{i+j=\ell}a_i(S)b_j(S).
\]
Apply subset M\"obius inversion to $c_\ell$. A term indexed by input subsets
$U,V$ remains at $S$ exactly when $U\cup V=S$. Selecting $\ell=|S|$ also
requires $|U|+|V|=|S|$, hence $U\cap V=\varnothing$. The selected coefficient
is \eqref{eq:conv}. In characteristic two, the subset M\"obius transform is
the same XOR butterfly as the zeta transform.

\subsection{Exact rank lanes in Python integers}
The implementation packs ranks into separated binary lanes. Put
\[
 \ell=\floor{\log_2(k+1)}+1,
 \qquad B=2^\ell>k+1,
 \qquad A(S)=\sum_{i=0}^k a_i(S) B^i.
\]
The initial arrays contain $a(S)B^{|S|}$ and $b(S)B^{|S|}$. XOR zeta
butterflies produce the packed arrays of rank parities. Each rank digit
is zero or one. In the ordinary integer product $A(S)B'(S)$, a coefficient
of any power of the lane base is at most $k+1$. Thus there is no carry
between lanes. Keeping the low bit of each lane computes rank-polynomial
multiplication modulo two exactly. The subsequent XOR M\"obius transform
and extraction of lane $|S|$ produce $h(S)$.

This is integer arithmetic throughout. There is no floating-point
convolution, randomized polynomial identity test, or probabilistic error.
The use of a lane base \emph{strictly} exceeding $k+1$ matters; treating an
ordinary integer product as carryless without this separation would be
incorrect.

\subsection{Projection, reconstruction, and complexity}
For a left monomial mask $U\ne0$, remove its least set bit $i$. If the
projection of $U\setminus\{i\}$ is already zero, or already contains the
component containing $i$, record zero. Otherwise adjoin that component.
This constructs a projection table of length $2^r$ with one recurrence per
entry. The right table is analogous. The output table uses
Lemma~\ref{lem:output}. These tables can be reused for a fixed topological
plan.

Let $\mathsf M(q)$ denote an upper bound for multiplying $q$-bit integers,
and let $W=(k+1)\ceil{\log_2(k+2)}$. A direct accounting of the packed
transforms gives the following statement.

\begin{theorem}[Dense composition bound]
\label{thm:bitcost}
A genus-zero plan with dimensions $(r,s,t,k)$ can be compiled and evaluated
exactly with bit cost bounded by
\begin{equation}
 \poly(r+s+t+k)\,(2^r+2^s+2^t)
 +O\!\left(2^k\bigl(kW+\mathsf M(W)\bigr)\right).
 \label{eq:bitcost}
\end{equation}
In particular, its bit cost is
$\poly(r+s+t+k)2^{\max\{r,s,t,k\}}$. Space has the same exponential parameter,
up to polynomial factors.
\end{theorem}
\begin{proof}
The two projection tables store one $O(k)$-bit component mask, or a failure
marker, per input subset. Building them costs a polynomial number of bit
operations per entry. Reading a packed input as bytes and accumulating
surviving terms takes at most polynomial overhead per coefficient position.
The output table is built by scanning the $k$ blocks for each of $2^t$
subsets, and reconstruction is one coefficient lookup per output subset.

The ranked arrays have $2^k$ entries of $O(W)$ bits. There are $O(k2^k)$
XOR butterflies and $2^k$ products of $O(W)$-bit integers. Lane extraction
and serialization add only linear passes with polynomial overhead.
Together these give \eqref{eq:bitcost}. Ordinary polynomial-time integer
multiplication suffices for the simplified bound; no special machine-word
assumption is needed.
\end{proof}

\begin{corollary}[Boundary-size bound]
\label{cor:width}
For composition between matching objects on $w=2m$ boundary points, the
new dense algorithm has bit cost $\poly(w)2^{w/2}$.
\end{corollary}
\begin{proof}
Each of the three circle spaces is formed from two matchings and has at
most $m$ circles. Each gluing component contains at least one left input
disc and at least one right input disc, so $k\leq\min(r,s)\leq m$.
Apply Theorem~\ref{thm:bitcost}. The crossing-free case is scalar arithmetic.
\end{proof}

The reference support-pair loop can enumerate $2^{r+s}$ pairs even before
its output arithmetic is counted. In a dense matching endomorphism algebra,
$r=s=t=k=m$; its pair enumeration is of order $4^m$, while the new dense
bit algorithm is $\poly(m)2^m$. This compares actual encoded arithmetic,
not a fictitious unit-cost operation on exponentially long integers.

\begin{remark}[Budgeted code versus the unbudgeted theorem]
The shipped adapter has an allocation ceiling and a sparse fallback.
When the ceiling excludes dense tables it deliberately uses the existing
exact algorithm. That capped adapter does not inherit
Corollary~\ref{cor:width} on arbitrarily large widths. The theorem concerns
the dense algorithm with sufficient memory; the budget policy concerns
practical safety. An entry-count cap is also not a byte-count cap.
\end{remark}

\section{Optimal information compression and explicit planar families}
\label{sec:optimal}
\subsection{The exact observational quotient}
\begin{definition}
For a fixed plan, two left inputs are observationally equivalent if
\[
 f\sim_L f'\quad\Longleftrightarrow\quad
 \Comp(f,g)=\Comp(f',g)\text{ for every right input }g.
\]
\end{definition}

\begin{theorem}[Universal-summary optimality]
\label{thm:optimal}
Assume that all components have genus zero and no extra dots, and that
every $L_j$ and $R_j$ is nonempty. Then
\[
 f\sim_L f'\quad\Longleftrightarrow\quad P_L(f)=P_L(f').
\]
There are exactly $2^{2^k}$ left observational equivalence classes.
Every linear summary sufficient for all right-input composition queries
has rank at least $2^k$, and the component quotient attains this rank.
Every fixed-length deterministic summary sufficient for those queries,
even a nonlinear one, requires at least $2^k$ bits. The same conclusions
hold with left and right exchanged.
\end{theorem}
\begin{proof}
The forward sufficiency of equal projections follows from
Theorem~\ref{thm:factor}. For necessity, let
$a=P_L(f-f')\ne0$. Both projections are surjective: choose one input
circle on each component and map the corresponding square-free monomials
onto the basis of $\R_k$.

Let $\lambda$ be the coefficient functional selecting the monomial with
\emph{all} output circles dotted. When there are no output circles this
means the scalar coefficient. By Lemma~\ref{lem:output},
\[
 \lambda\Omega(h)=[y^{[k]}]h.
\]
The bilinear pairing $(a,b)\mapsto[y^{[k]}]ab$ on $\R_k$ is nondegenerate:
$y^S$ and $y^{[k]\setminus S}$ are dual basis vectors. Select $S$ with
$a_S=1$ and take $b=y^{[k]\setminus S}$. Surjectivity supplies a right
input $g$ with $P_R(g)=b$. Then
\[
 \lambda\bigl(\Comp(f,g)-\Comp(f',g)\bigr)=1,
\]
so the two left inputs are distinguishable.

The quotient is therefore the $2^k$-dimensional space $\R_k$ itself, with
$2^{2^k}$ elements. A sufficient linear summary has kernel contained in
$\ker P_L$, giving the rank bound. An arbitrary sufficient deterministic
summary must have distinct values on all these classes, giving the
fixed-length information bound. The symmetric argument handles the right.
\end{proof}

This is not a lower bound for all methods of deciding whether a knot is
trivial. It is a lower bound for a precisely specified universal
composition-summary interface on arbitrary coefficient vectors. It also
does not rule out shorter descriptions of special input families.

\begin{corollary}[Forced-dot refinement]
\label{cor:forced}
Under the same nonempty-input assumptions, suppose the extra-dot set is
$E$ and no component has two extra dots. Let $q_E$ set $y_j=0$ for $j\in E$.
The exact left observational quotient is $q_E P_L$, of dimension
$2^{k-|E|}$.
\end{corollary}
\begin{proof}
For any $a,b\in\R_k$,
$y^Eab=y^E q_E(a)q_E(b)$. The top-output coefficient now pairs the residual
algebra on $[k]\setminus E$ by its top monomial. The proof of
Theorem~\ref{thm:optimal} applies to that residual algebra.
\end{proof}

The code handles extra dots correctly, but does not yet use this smaller
convolution dimension. Extra dots are absent in ordinary matching
composition; the refinement is relevant to more general compiled maps.

\subsection{Planar zipper triples}
An arithmetic benchmark should not hide an impossible topology. We now
give an explicit family of genuine planar matching triples.

\begin{proposition}[Zipper family]
\label{prop:zipper}
For any positive integers $d,q$, there are three noncrossing matchings
$a,b,c$ on $2q(2d-1)$ commonly ordered boundary points such that the
composition plan has
\[
 r=s=qd,\qquad t=k=q,
\]
with $q$ genus-zero components. Each component has $d$ left input circles,
$d$ right input circles, and one output circle.
\end{proposition}
\begin{proof}
Construct one block with $m=2d-1$ arcs. Let $b$ pair consecutive points,
so its arcs are $(0,1),(2,3),\ldots,(2m-2,2m-1)$.
To obtain $a$, replace the consecutive arc pairs with starting indices
$j=0,2,\ldots,2d-4$ by
\[
 (2j,2j+3),\quad(2j+1,2j+2).
\]
To obtain $c$, make the same replacements starting at
$j=1,3,\ldots,2d-3$. Each replacement lies in its own four-point interval,
so the resulting matchings are noncrossing.

Each replacement merges two circles of $a\cup\overline b$, or of
$b\cup\overline c$, and the replacements are disjoint within each
matching. Both circle counts are therefore $m-(d-1)=d$.
Make a bipartite graph whose vertices are these left and right circles
and whose edges are the $m$ arcs of $b$. The alternating pairings make
this graph a path on $2d$ vertices. Gluing the corresponding $2d$ discs
along its $2d-1$ edges produces a connected disc. Hence the gluing has one
component, genus zero, and one output boundary circle.
Take $q$ disjoint consecutive copies of this construction.
\end{proof}

For fixed $q$, the component algebra has dimension $2^q$ even though the
input spaces each have dimension $2^{qd}$. The required input scan is still
present: compression does not make reading a dense input free. The
benchmark with $(d,q)=(5,2)$ uses this construction on 36 boundary points.
The exact matchings and masks are stored in \texttt{results/examples.json},
and 24 zipper triples are checked by the tests.

A useful consistency identity follows from Euler characteristic. If a
matching composition has $m$ glued arcs, $k$ components, $t$ output
circles, and total genus $G$, then
\begin{equation}
 r+s-m=2k-2G-t.
 \label{eq:euler}
\end{equation}
The left side counts input discs minus glued intervals. In the genus-zero
case this gives $r+s+t-m=2k$. The zipper tests check this identity as well
as the explicit component masks.

\section{Integration and the remaining multiplicity problem}
\label{sec:integration}
\subsection{What the adapter changes}
\texttt{ComponentAlgebra} wraps a new scanner's existing \texttt{Planar}
instance. It delegates geometry, circle numbering, stage setup, transfer,
and existing shape/result caches. A stage-local cache stores compiled
component quotients. Only sufficiently dense composition calls use the
new kernel. Small support products and excluded allocations use the
original composition. No global monkey-patch is used.

The default policy first requires at least 256 support pairs. It then
compares that count with a rough table-and-transform work estimate. This
is a transparent heuristic, not a theorem predicting Python wall time.
Compiling a table also costs time and memory; amortization depends on
actual plan reuse. Source-level identity shortcuts and complete result
cache hits occur before this numerical kernel and should remain there.

The check callback propagates the scanner's existing budget or race
exception. A dense-allocation failure means ``use the exact old kernel,''
not ``this is a knot.'' The wrapper is owned by one scanner, and hence by
one race worker if the caller installs it there. Cross-process installation
and the production command-line interface have not been changed.

\subsection{A parameter ledger for complete scans}
Let $n$ be the number of scanned crossings, $w$ the maximum number of
boundary endpoints, and $Q$ the maximum number of explicit objects present
at a stage \emph{before} cancellation. Counting only the objects left after
elimination would miss the transfer peak.

\begin{proposition}[Dense-primitive scanner envelope]
\label{prop:scanner}
A conventional explicit scanning-and-cancellation algorithm can be
implemented with bit cost
\begin{equation}
 n\,\poly(n,Q,w)\,2^{O(w)}.
 \label{eq:envelope}
\end{equation}
If composition uses Theorem~\ref{thm:bitcost} and linear transfer maps use
dense coefficient reconstruction as described below, an implementation has
the sharper envelope
\begin{equation}
 n Q^3\,\poly(n,w,\log(Q+1))\,2^{w/2}.
 \label{eq:sharper}
\end{equation}
These are parameterized implementation envelopes, not uniform bounds on
$Q$ or $w$ for all knot diagrams.
\end{proposition}
\begin{proof}
A stage has at most $Q$ cancellations. Cancelling a unit creates at most
$Q^2$ source-target update pairs, each requiring a bounded number of
compositions. For this existence bound, a pivot can be selected by scanning at most
$Q^2$ entries per cancellation; an optimal Markowitz choice is not required.
Explicit arrays give deterministic access with polynomial bit overhead
and logarithmic identifier sizes. Geometry on at most $w$ endpoints is
polynomial. Thus $O(nQ^3)$ composition updates suffice as an envelope.

Adding one crossing has a bounded number of smoothing and delooping
choices per old object, and at most polynomially many entries in $Q$.
For a fixed linear genus-zero transfer map, aggregate input monomials by
component masks in a balanced search tree with at most $2^r$ keys. For each of the
$2^t$ output subsets, Lemma~\ref{lem:output} identifies at most one key
to query. With no second input to convolve, this takes
$\poly(w)(2^r+2^t)$ bit operations. Crossing transfer has only a bounded
number of such maps; untouched circles are included as one-input,
one-output components. Here $r,t\leq w/2$ at the adjacent frontiers.
Together with Corollary~\ref{cor:width} this proves \eqref{eq:sharper}.

Even direct support expansions and pair evaluation have only $2^{O(w)}$
terms and $2^{O(w)}$-bit intermediate morphisms. Combining these with the
polynomial number of object updates gives \eqref{eq:envelope}.
\end{proof}

The shipped adapter leaves the existing transfer code unchanged. Therefore
\eqref{eq:sharper} is an envelope for the stated dense-primitive
implementation, not a measured or fully implemented bound for that adapter.
The weaker exponential-in-$w$ envelope is sufficient for the following
criterion.

\begin{corollary}[A sufficient scanning criterion]
\label{cor:scanqp}
Suppose a family has polynomial-size input encodings and computable scan
schedules whose construction costs $n^{O(\log n)}$. If the schedules satisfy
\[
 w+\log_2(Q+1)=O\bigl((\log n)^2\bigr),
\]
then the parameterized implementation in \eqref{eq:envelope} computes its
scans in $n^{O(\log n)}$ time.
\end{corollary}
\begin{proof}
Take binary logarithms of \eqref{eq:envelope}. A fixed polynomial in
$n,Q,w$ contributes $O(\log n+\log(Q+1)+\log(w+1))$, and the exponential
factor contributes $O(w)$. The assumed schedule cost has the same
quasi-polynomial upper bound.
\end{proof}

Let $M$ be the number of distinct matching types at a stage, and let $\mu$
be the maximum number of copies of any one matching, summed over all
homological degrees. Then $Q\leq M\mu$. If all matchings are noncrossing in
one certified cyclic order on the boundary of a disc,
$M\leq \mathrm{Cat}_{w/2}\leq2^w$. Without that geometric hypothesis, the
safe purely combinatorial bound for pairings is
$M\leq(w-1)!!$, giving
\[
 \log Q=O(w\log(w+2))+\log(\mu+1).
\]
A frontier-size argument must identify which setting it uses. We have not
audited all repository scan frontiers to upgrade an observed matching
count to a universal Catalan-bound theorem.

\subsection{Fixed boundary does not control explicit multiplicity}
\label{sec:multiplicity}
\begin{proposition}[Connected-sum multiplicity obstruction]
\label{prop:mult}
Choose a nontrivial knot $K_0$ and write
$\beta=\dim_{\F}\widetilde\Kh(K_0;\F)>1$.
There is a family of $(1,1)$ tangles $T_q$ with two boundary endpoints,
whose reduced closures have Khovanov dimension $\beta^q$.
Any explicit complex of copies of the single boundary matching that is
chain-homotopy equivalent to the tangle complex of $T_q$ must contain at
least $\beta^q$ objects in total.
\end{proposition}
\begin{proof}
Cut a marked strand in the connected sum of $q$ copies of $K_0$ to obtain
$T_q$. At the cube-of-resolutions level, the complex of a connected sum is
the tensor product over the marked-circle algebra $A$. Each resolution
is free over that marked factor. Applying the reduced base change
$A\to A/(x)=\F$ identifies the reduced connected-sum complex with the
tensor product over $\F$ of the reduced factor complexes. The field
K\"unneth formula gives reduced homology dimension $\beta^q$.

Under reduced closure, each copy of the one-arc boundary matching yields
one basis vector. Applying this additive functor to a chain homotopy
equivalence preserves that equivalence. Homology dimension cannot exceed
total chain-group dimension, so at least $\beta^q$ objects are necessary.
Finally, $\beta>1$ follows from the rational reduced unknot-detection
theorem \cite{km} and the universal-coefficient inequality
$\dim_{\F}\widetilde\Kh\geq\dim_{\mathbb Q}\widetilde\Kh$.
\end{proof}

The crossing count of this family is linear in $q$ for a fixed diagram of
$K_0$. Thus no polynomial multiplicity bound can depend on the two-point
boundary alone. This is not a hardness proof for unknot recognition: the
family consists of visibly composite nontrivial knots, and the repository's
existing connected-sum stage is designed to avoid such redundant full
computations. The point is narrower and useful: an unrestricted proof may
not count boundary matchings and silently discard their multiplicities.

\section{Sparse hierarchy restart entropy}
\label{sec:entropy}
\subsection{The source of the lexicographic profile}
Lackenby's 2026 preprint uses integer pattern-complexity
$\chi_p(S)=-4\chi(S)+|S\cap P|$. In Section~9, nonnegative layer
complexities are padded to length $L$; between successive construction-loop
exits, a first changed coordinate decreases and its suffix is rebuilt.
The resulting bound is $L(g+1)^L$ iterations when layer values are at most
$g$. This counts iterations, not their bit costs \cite[Section~9]{lackenby}.

The refinement below is a theorem about that restart pattern. Its new
hypothesis is a bound on the number of positive coordinates in every
relevant profile. No assertion that all knot exteriors satisfy this
hypothesis is made.

\begin{definition}[Sparse profile space]
For nonnegative integers $L,g,s$, set
\[
 \mathcal V(L,g,s)=\{x\in\{0,1,\ldots,g\}^L:
       |\{i:x_i>0\}|\leq s\}.
\]
A restart trace is a sequence of these vectors in strictly decreasing
lexicographic order. One may replace all coordinates after the first
changed coordinate arbitrarily, subject to membership in $\mathcal V$.
\end{definition}

The actual positions of zero-complexity layers are not erased. They can
carry topology and affect where positive layers occur. This is why the
count below contains $\binom Lj$, rather than merely $g^s$.

\begin{theorem}[Sharp sparse restart bound]
\label{thm:entropy}
The profile space has cardinality
\begin{equation}
 N(L,g,s)=\sum_{j=0}^{\min\{s,L\}}\binom Lj g^j.
 \label{eq:N}
\end{equation}
A strictly descending trace has at most $N(L,g,s)$ states and at most
$N(L,g,s)-1$ restart transitions. If an epoch between restarts uses at most
$C(L+1)$ primitive calls, the entire trace and its final epoch use at most
$C(L+1)N(L,g,s)$ calls. The state bound is sharp under the stated abstract
restart rules.
\end{theorem}
\begin{proof}
To choose a vector with exactly $j$ positive coordinates, choose their
locations in $\binom Lj$ ways and assign each of them one of $g$ positive
values. Summing gives \eqref{eq:N}, including the all-zero vector.
Strict descent cannot revisit a vector, so it has the stated length.
Multiplying the number of epochs by their assumed call bound proves the
algorithmic count.

For sharpness, list every vector of $\mathcal V$ in decreasing
lexicographic order. Between consecutive vectors, the first differing
coordinate decreases; all earlier coordinates agree. The later coordinates
are a permissible replacement suffix. Hence this full list is a legal
abstract restart trace with exactly $N$ states.
\end{proof}

\begin{corollary}[Polynomial total depth, logarithmic positive depth]
\label{cor:sparseqp}
For $n\geq2$, if $L,g\leq n^a$ and $s\leq b\log_2 n$ for fixed constants
$a,b$, then
\[
 N(L,g,s)=n^{O(\log n)}.
\]
The same is true after multiplying by a polynomial bound on the number of
primitive calls per epoch.
\end{corollary}
\begin{proof}
For $Lg\geq1$,
\[
 N(L,g,s)\leq(s+1)\max\{1,Lg\}^{s}.
\]
Taking logarithms gives
$\log_2 N\leq\log_2(s+1)+s(\log_2L+\log_2g)=O((\log n)^2)$.
When $L=0$ or $g=0$, the space has one vector.
\end{proof}

A mere polynomial bound on $L$ is insufficient without another restriction:
with $g=1$ and unrestricted support, reverse binary counting gives $2^L$
states. Theorem~\ref{thm:entropy} is therefore not a repackaging of
``polynomial depth implies quasi-polynomial time.'' It identifies precisely
which additional combinatorial resource would suffice.

\subsection{An exact executable rank potential}
Let $\rho(x)$ be the number of members of $\mathcal V(L,g,s)$ strictly
smaller than $x$ in increasing lexicographic order. While scanning a prefix
of $x$, suppose $r$ nonzero positions remain available and there are $m$
coordinates after the current one. If the current digit $d$ is positive,
its contribution to the rank is
\begin{equation}
 N(m,g,r)+(d-1)N(m,g,r-1).
 \label{eq:rank}
\end{equation}
The first term counts the smaller choice zero. The second counts smaller
positive digits. Continue with remaining budget $r-1$. A zero current digit
contributes nothing. This proves the rank recurrence used by
\texttt{lex\_rank}. Its inverse compares the rank with the zero-digit block
size and then uses division by the positive-digit block size.

The potential is an exact nonnegative integer, decreases at every restart,
and is initially at most $N-1$. It uses only $\ceil{\log_2N}$ bits.
Evaluation and inverse evaluation use binomial coefficients and integer
arithmetic polynomial in their explicit dimensions and bit lengths. The
trace tool is not an oracle for topology: it checks a supplied finite trace
and cannot prove that an unspecified algorithm will always preserve the
sparsity bound.

For $L=4$, $g=2$, and $s=2$, the count is $1+8+24=33$. The package supplies
the complete reverse-ranked trace of these 33 vectors. It also checks
23,460 rank/unrank identities and 140 complete countdowns.

\subsection{Other entropy budgets}
\begin{proposition}[Heterogeneous and additive budgets]
\label{prop:otherbudgets}
If coordinate $i$ has a fixed cap $g_i$, the unrestricted state count is
$\prod_i(g_i+1)$. With support at most $s$, it is
\[
 \sum_{j=0}^{s}[z^j]\prod_{i=1}^L(1+g_i z).
\]
If instead nonnegative coordinates satisfy $\sum_i x_i\leq G$, their
state count is $\binom{L+G}{G}$, before imposing any further coordinate caps.
\end{proposition}
\begin{proof}
The first formula is the product rule. In the generating function, the
constant term at a position selects zero, and $g_i z$ selects one of its
positive values. For the additive budget, adjoin a nonnegative slack
coordinate to make the total equal to $G$ and apply stars and bars.
\end{proof}

Thus another possible target is total charge $G=O(\log n)$ with polynomial
$L$. Fixed-coordinate caps have entropy $\sum_i\log_2(g_i+1)$.
These alternatives can be useful when a maximum cap $g$ is a poor summary.
The caps must be valid uniformly over the trace, not chosen retrospectively
only for the currently observed vector.

\subsection{What would complete a hierarchy-based quasi-polynomial proof}
\begin{theorem}[Conditional hierarchy criterion]
\label{thm:conditional}
Suppose an exact hierarchy-based recognition algorithm on inputs of size
$n$ has all of the following properties. Its restart profiles lie in
$\mathcal V(L,g,s)$ and strictly decrease; $L,g$ are polynomially bounded
and $s=O(\log n)$; at most polynomially many primitive calls occur in an
epoch; and every primitive has a uniform $n^{O(\log n)}$ bit-time bound,
including construction and maintenance of its encoded state. Assume its
initialization has the same bound and its terminal decisions are sound and
complete. Then the algorithm recognizes the unknot in $n^{O(\log n)}$ time.
\end{theorem}
\begin{proof}
Corollary~\ref{cor:sparseqp} bounds the number of epochs by
$2^{O((\log n)^2)}$. Multiplying this by polynomially many calls per epoch
and by $2^{O((\log n)^2)}$ work per call preserves that form. Add the
initialization cost and use the assumed decision correctness.
\end{proof}

This is a sufficient-condition theorem, not a completed topological
algorithm. The difficult missing statement is a \emph{constructive},
uniform sparsity or comparable entropy bound for all states the algorithm
actually reaches, including rebuilt suffixes and failed choices.
An efficiently describable successful certificate alone does not supply
that statement.

The source already proves important compressed ingredients: Section~10
controls refined hierarchy length linearly in initial $q$-complexity;
Theorem~14.1 gives polynomial-size verifiable essential hierarchies;
Proposition~14.2 gives polynomial-time cut-and-simplify operations under its
hypotheses, measured in handle count and logarithmic extended surface
weight \cite{lackenby}. These should be reused, not described as absent.
What remains for our criterion is uniform efficient construction, valid
hypotheses throughout the run, and low restart entropy. Zero potential
increments do not make geometric operations free.

\section{Validation and performance experiments}
\label{sec:experiments}
\subsection{What was executed}
The artifact is standard-library Python, tested here on Python 3.13.5,
GCC 14.2.0, Linux x86\_64. The final recorded run completed 15 test methods
with zero failures and zero errors in 7.07 seconds. This is a numerical
kernel and abstract-potential test suite, not the repository's full suite.

\begin{center}
\begin{tabular}{lr}
\toprule
Validation category & Recorded checks\\
\midrule
Subset-convolution comparisons and square identities & 65,980\\
Matching-composition equalities & 460,564\\
General plans: counits, extra dots, and splitting & 24,000\\
Associativity identities & 1,000\\
Adaptive-adapter comparisons & 31\\
Explicit zipper triples & 24\\
Sparse-profile rank/unrank identities & 23,460\\
Complete sparse countdown traces & 140\\
\bottomrule
\end{tabular}
\end{center}

All coefficient pairs are checked for all planar matching triples through
six boundary points whenever their plan has genus zero. Larger randomized
triples cover eight through fourteen endpoints. The 2,535 sampled or
exhaustively enumerated triples include 590 positive-genus plans, which
are counted separately rather than as successful dense-kernel evaluations.
The general-plan tests cover output counits, zero maps, multiple extra dots,
and empty universes. Invalid masks, out-of-basis coefficients, allocation
limits, stage resets, and cooperative budget exceptions also have tests.

The reference is an excerpt of the repository's actual numerical
algorithms, not a second full topological implementation. Equality with it
is useful regression evidence but could share a source convention error.
The monomial proof of Theorem~\ref{thm:factor}, the independent
subset-convolution checks, and the associativity checks supply different
kinds of evidence. None substitutes for a full scanner regression suite or
a formal proof assistant verification.

\subsection{Paired numerical-kernel measurements}
The benchmark uses 11 rounds, seed 20261007, and alternating ABAB/BABA
measurement order. Each round checks output equality and records two old
and two new timings. A/A ratios are retained. Topology is precompiled for
both implementations; component-table compilation is timed separately.
The reference numerical monomial memo is fresh on each invocation. The
new method has no numerical-result cache. This measures a particular
cold-memo numerical workload, not the amortized behavior of a complete
scanner with all its caches.

\begin{center}
\small
\begin{tabular}{lrrrr}
\toprule
Case & Old (ms) & New (ms) & Paired speed ratio & Compile (ms)\\
\midrule
Dense endomorphism, $k=4$  & 0.0322 & 0.0206 & 1.56 & 0.0444\\
Dense endomorphism, $k=6$  & 0.5877 & 0.1162 & 5.03 & 0.0908\\
Dense endomorphism, $k=8$  & 9.8212 & 0.4247 & 22.89 & 0.2890\\
Dense endomorphism, $k=10$ & 179.5399 & 2.1522 & 86.72 & 1.3189\\
Sparse endomorphism, $k=10$ & 0.0126 & 1.2923 & 0.00972 & 1.3195\\
Zipper, $(r,s,t,k)=(10,10,2,2)$ & 69.1873 & 0.1742 & 401.98 & 0.2521\\
\bottomrule
\end{tabular}
\end{center}

Old and new columns are separate medians; the speed ratio is the inverse
of the median paired new/old ratio, so it need not equal their quotient.
The sparse case has four terms in each operand and deliberately forces the
dense kernel. It is approximately 103 times slower by the paired median.
The adapter's default 256-pair gate avoids that particular substitution.

The observed paired new/old ranges for dense $k=6,8,10$ are respectively
$[0.142,0.234]$, $[0.0336,0.0540]$, and $[0.0105,0.0128]$.
The zipper range is $[0.00201,0.00267]$.
These are ranges of observed rounds, not confidence intervals.
The $k=4$ range extends above one, so its small median gain is not a
robust universal speed claim. The A/A controls and every raw timing are
available in \texttt{results/benchmark.json}.

These are genuine matching-algebra workloads, including the explicit
planar zipper family, but they are \emph{not measured recognition speedups
on knot diagrams}. Existing pipeline filters may settle an input before
Khovanov computation; existing caches may remove numerical work; and many
scanner operands may remain sparse. Full-input paired timing, warm-memo
behavior, peak memory, and the frequency of compressed calls are necessary
before enabling the adapter by default.

\subsection{Reproduction and repository validation}
From the package root:
\begin{lstlisting}
python tools/run_tests.py
python tools/benchmark.py
pdflatex -interaction=nonstopmode -halt-on-error article.tex
pdflatex -interaction=nonstopmode -halt-on-error article.tex
\end{lstlisting}
The \texttt{Makefile} provides the corresponding targets. Timing values
will vary by machine and run; structural counts should be reproducible.
The article's table records the delivered run, rather than automatically
changing when somebody reruns the benchmark.

\texttt{integration/check\_repository.py} is a read-only local-checkout
harness. It reports audited source hashes, compares the kernel against the
actual \texttt{Planar}, and optionally compares direct closed-PD scanner
runs with $d^2=0$ checks. It was supplied but not executed against a complete
checkout in this environment. No Rust port, full repository suite, or
production recognizer benchmark is claimed. The integration guide states
these limits and the required pre-adoption checks explicitly.

\section{Further research questions and completion criteria}
\label{sec:research}
\subsection{Production relevance and adaptive arithmetic}
\begin{question}
On the current hard-input corpus, how often does composition reach the
dense regime after the existing filters, cancellation, identity shortcuts,
and caches have run?
\end{question}
A useful result is a paired full-scanner study with compressed-call counts,
operand densities, component counts, table-reuse counts, and peak memory.
The decision criterion should be improvement in completed recognition or
scan time, not merely a large numerical-kernel ratio. Preserve cases with
zero compressed calls in the report. An implementation milestone is default
enablement only after complete Python cross-checks and a separate Rust
port validation.

\begin{question}
Can a dispatcher choose between sparse pair evaluation, component
projection, and packed convolution with bounded overhead relative to the
best available arithmetic strategy?
\end{question}
The current support threshold is heuristic. A next version should measure
warm-memo reuse and compilation amortization and should include the
forced-dot refinement of Corollary~\ref{cor:forced}. A theoretical target
is a transparent online work bound under a declared primitive-cost model;
a practical target is paired speed with no sparse-input regression beyond
a specified tolerance. Dense-linear transfer reconstruction is another
well-defined implementation task needed for the sharper scanner envelope.

\subsection{Representations that do not enumerate multiplicity}
\begin{question}
For which tangle families does a tensor, circuit, or module presentation
retain the cancellation information needed for an unknot decision while
avoiding explicit copies of every matching object?
\end{question}
Proposition~\ref{prop:mult} excludes a width-only bound for the explicit
representation, not for every representation. A successful theorem must
specify how composition, testing units, cancellation, and the final decision
operate on the compressed presentation. Short descriptions of chain groups
alone do not bound the cost of deciding their homology. First targets could
be families closed under a specified tangle composition with proved bounds
on description size and update cost.

\begin{question}
Can one maintain a decision-sufficient certificate instead of computing the
entire reduced Khovanov rank on difficult nontrivial inputs?
\end{question}
The universal-summary theorem concerns all opposite-operand composition
queries, a much stronger demand than one final decision. A useful advance
would identify a smaller query class justified by the scanner and prove
that its retained information is sound under every continuation. A
shortcut must not interpret failure to discover a second homology class as
proof of rank one. Existing connected-sum reduction should remain part of
this investigation, rather than being counted as a new result.

\subsection{Low-entropy hierarchy construction}
\begin{question}
Is there a constructive hierarchy strategy for knot exteriors for which
every restart profile has at most $C\log n$ positive pattern-complexity
layers, with polynomial coordinate caps and total padded length?
\end{question}
The required bound must survive unsuccessful surface choices, suffix
reconstruction, and simplification. A theorem only about a best final
certificate is insufficient. Theorem~\ref{thm:conditional} gives a precise
completion criterion: prove this uniform sparsity and the required
primitive bit costs, or replace sparsity by another profile entropy bound
of order $(\log n)^2$.

\begin{question}
Can zero-complexity layers be represented by composable geometric records
whose replay, transport of patterns, and invalidation after a restart have
polynomial total cost?
\end{question}
Zero pattern-complexity does not mean a layer is an identity operation.
The representation must preserve boundary attachments and certificates
that local replacements induce the intended manifold operations. A
promising target is a persistent directed acyclic record of cuts and
parallelity reductions with explicit size, composition, and verification
bounds. A necessary audit separates constructing a useful surface from
cutting along a surface that has already been supplied.

\begin{question}
Which large abstract restart countdowns are realizable by actual
hierarchies and a fixed deterministic choice strategy?
\end{question}
The sharpness construction in Theorem~\ref{thm:entropy} is combinatorial,
not a topological lower bound. Explicit diagram families, including prime
and iterated satellite families, could distinguish genuine geometric
constraints from artifacts of an overpermissive state model. Conversely,
a proved family with many positive layers would refute the naive uniform
sparsity conjecture for that particular strategy and motivate heterogeneous
or additive charge budgets.

\subsection{Verification and scan scheduling}
\begin{question}
Can the component factorization, its observational quotient, and the
sparse-rank potential be formalized with extraction of the executable
arithmetic, then connected to verified geometric primitives?
\end{question}
The finite-set algebra offers a relatively small first formalization:
quotient-ring maps, the top-coefficient pairing, and ranked M\"obius
inversion. The restart theorem requires finite lexicographic order and
exact counting. The geometric interface is a separate proof obligation:
compiled masks must describe the intended surface, and recorded hierarchy
transitions must preserve their topological invariants. Regression tests
should continue to target those interfaces.

A related scheduling task is to predict work from both boundary width and
object multiplicity. The profile $(w,Q,k)$ and observed morphism density
are more informative than width alone. Any new schedule search must count
its own running time and must be compared with the repository's existing
order selection and racing rather than treated as free preprocessing.

\section{Conclusion}
The immediate implementable advance is an exact dense-arithmetic kernel
for the existing scanning backend. Its governing variable is the number
of connected gluing components, and its composition cost on $w$ endpoints
is $\poly(w)2^{w/2}$ in the dense bit model. The exact observational quotient
and the planar zipper family explain both why the compression is valid and
where substantial reductions can occur. Sparse workloads, memory caps,
transfer costs, and multiplicity remain visible in the analysis.

For the global asymptotic objective, sparse restart entropy is a distinct
and concrete target. Polynomial hierarchy length is compatible with
quasi-polynomial restart counts when only logarithmically many coordinates
are positive. What must still be proved is a constructive topological
bound of that kind, together with complete primitive cost accounting.
The supplied proofs, code, negative cases, test records, and integration
harness are intended to make those next steps experimentally accessible
without overstating what has already been established.

\appendix
\section{Artifact map and audit boundaries}
\begin{center}
\small
\begin{tabular}{p{0.39\linewidth}p{0.53\linewidth}}
\toprule
Path & Role\\
\midrule
\texttt{article.tex}, \texttt{article.pdf} & Self-contained article source and compiled text.\\
\texttt{code/component\_kernel.py} & Exact quotient, convolution, reconstruction, and opt-in adapter.\\
\texttt{code/restart\_entropy.py} & Counts, lexicographic ranks, inverse ranks, finite-trace audit.\\
\texttt{tests/reference\_planar.py} & Numerical excerpt from the audited repository source; not a full recognizer.\\
\texttt{tests/fixtures.py} & Explicit planar zipper construction.\\
\texttt{tests/test\_*.py} & Regression and structural checks.\\
\texttt{tools/run\_tests.py}, \texttt{tools/benchmark.py} & Reproducible validation and paired kernel measurements.\\
\texttt{results/} & Raw JSON, console records, worked examples, and finite traces.\\
\texttt{integration/} & Opt-in usage and unexecuted local-checkout validation harness.\\
\texttt{PROVENANCE.json}, \texttt{CLAIMS.md} & Exact source identifiers and claim/limitation ledger.\\
\bottomrule
\end{tabular}
\end{center}

No third-party lecture slides or research-paper PDFs are redistributed.
The source excerpt comes from the repository's MIT-0 material. The package
contains no automatic modification of the production repository.
The root README and claim ledger distinguish delivered execution results
from suggested future validation commands.

\begin{thebibliography}{9}
\bibitem{repo}
V. Reshetnikov and contributors,
\emph{ProveIt: Topology/UnknotRecognition}, source audited 7 October 2026.
\url{https://github.com/VladimirReshetnikov/ProveIt/tree/main/Topology/UnknotRecognition}.
Exact reviewed blob identifiers are recorded in \texttt{PROVENANCE.json}.

\bibitem{bn-tangles}
D. Bar-Natan,
\emph{Khovanov's homology for tangles and cobordisms},
arXiv:math/0410495.
\url{https://arxiv.org/abs/math/0410495}.

\bibitem{bn-fast}
D. Bar-Natan,
\emph{Fast Khovanov Homology Computations},
arXiv:math/0606318.
\url{https://arxiv.org/abs/math/0606318}.

\bibitem{bhkk}
A. Bj\"orklund, T. Husfeldt, P. Kaski, and M. Koivisto,
\emph{Fourier meets M\"obius: fast subset convolution},
arXiv:cs/0611101 (2006), Proceedings of STOC 2007.
\url{https://arxiv.org/abs/cs/0611101}.

\bibitem{km}
P. B. Kronheimer and T. S. Mrowka,
\emph{Khovanov homology is an unknot-detector},
arXiv:1005.4346;
Publications Math\'ematiques de l'IH\'ES \textbf{113} (2011), 97--208.
\url{https://arxiv.org/abs/1005.4346}.

\bibitem{lackenby}
M. Lackenby,
\emph{Incompressible surfaces, hierarchies and unknot recognition},
arXiv:2607.23350v1, 25 July 2026.
\url{https://people.maths.ox.ac.uk/lackenby/algorithm-incompressible-250726.pdf}.
\end{thebibliography}
\end{document}
